% page 500 du fac-simile (image p-499 du scan)
% bloc 162.6 x 224.7 mm ; marge gauche 8.0 mm
\pgc{-12.700mm}{0.000mm}{
\l{492}
\l{}
\l{\sou{renversar}. (\sou{trans.}) Faligar ulo, ulu, adsur sulo - generale}
\l{\cel{3}per shoko, pulso subita e ne-expektita da la pulsito.}
\l{}
\l{}
\l{\sou{reostato}. (\sou{elektr.)}\cel{2}Aparato qua inkluzas rezistivo elektra-}
\l{\cel{3}la, regulebla segun-vole, e qua, interpozite en cirkuito,}
\l{\cel{3}posibligas la variigo di la elektrofluo. - DEFIRS.}
\l{}
\l{\sou{reotaktismo}. (\sou{biol}.)Karakterizivo di ula organismi : li su}
\l{\cel{3}movas inverse di la aquo-fluo en qua li esas imersita.}
\l{}
\l{\sou{reotropismo}.\cel{2}(\sou{biol.}) Reotaktismo qua su manifestas che enti}
\l{\cel{3}komplexa (metafiti o metazoi).}
\l{}
\l{\sou{reparar}. (trans.)\cel{2}Igar lo necesa por igar(lu)aden lua stando}
\l{\cel{3}antea, pos lua domajiteso. - DEFIRS.}
\l{}
\l{\sou{repartisar}. (\sou{trans., inter)}. Partigar (ulo) ed atribuar a sin-}
\l{\cel{3}glu la parto qua apartenas a lu. - DEFIRS.}
\l{}
\l{\sou{repastar}. (\sou{netrans.,}\cel{1}ek).Manjar ye kloki determinita (dejunar,}
\l{\cel{3}dinear, supear ; matino-frue : dejunetar). - EFI.}
\l{}
\l{\sou{repentar}. (\sou{netrans., pri)}.\cel{2}Regretar sua kulpo, dezirante}
\l{\cel{3}emendar o kompensar lu, e ne plus agar lu. - EFIS.}
\l{}
\l{\sou{reperar}. (\sou{trans}.) I. Determinar la situeso topografiala pre-}
\l{\cel{3}ciza ed exakta. - II. (\sou{teknol.)}\cel{2}Markizar du bloki, du or-}
\l{\cel{3}gani, qui devas funcionar ensemble, por posibligar la sameso}
\l{\cel{3}absoluta di lia situeso rispektiva lor la ri-munto. -F.}
\l{}
\l{\sou{reperkutar}. (\sou{trans.}) Produktar retro-shoko. - (La korpo shokata,}
\l{\cel{3}shokas sua-foye la korpo quan lu kontaktas o preske kontak-}
\l{\cel{3}tas. - EFIS.}
\l{}
\l{\sou{repertorio}. I.Sorto de tabelo, en qua la tituli, la nomi, e c.}
\l{\cel{3}esas klasifikita segun ordino qua igas posibla trovar li}
\l{\cel{3}facile. - II. Listo de la teatraji qui permanas en latea-}
\l{\cel{3}tro e qui forsan ri-pleesos. - DEFIRS.}
\l{}
\l{\sou{repetar}. (\sou{trans.})\cel{2}Dicar plura-foye, parole o skribe, to quon}
\l{\cel{3}onu, od altru, ja dicis. - DEFIRS.}
\l{}
\l{\sou{repleta}. (\sou{pri vivanto}).Di qua la formo esas kelke ronda, la}
\l{\cel{3}karno ferma, la pelo pasable tensita. - EFS.}
\l{}
\l{\sou{replikar}. (\sou{trans., ad, pri)}. I. (a) Respondar a to quon ulu}
\l{\cel{3}respondis. (b)(\sou{yuro-cienco}) Respondar a la respondo da la}
\l{\cel{3}adverso. - II. Respondar a to quo, aspekte, ne postulas}
\l{\cel{3}respondo. - DEFIRS.}
}
